\section*{Введение}
\addcontentsline{toc}{section}{Введение}

\textbf{Задание лабораторной работы~1.}
\begin{enumerate}
    \item Сформировать случайным образом связный ациклический ориентированный граф
    в соответствии с заданным распределением (параметры распределения задаются как
    константы и подбираются экспериментально) с введённым пользователем количеством
    вершин. Распределение брать из справочника Вадзинского: реализовать генерирующую
    формулу распределения Вейбулла. Генерация происходит по степеням вершин.

    \item Посчитать экцентриситеты вершин, выделить центр графа и диаметральные вершины.

    \item Реализовать метод Шимбелла для сгенерированной с помощью заданного распределения
    весовой матрицы (пользователь вводит количество рёбер), в ответе представить минимальный
    и/или максимальный путь (в виде матрицы). При генерации необходимо учесть генерацию
    только положительных значений, только отрицательных значений и смешанную (на выбор
    пользователя).

    \item Определить возможность построения маршрута от одной заданной точки до другой
    (вершины вводит пользователь) и указывать количество таковых маршрутов.

    \item Реализовать меню для управления всеми пунктами.
\end{enumerate}

\textbf{Задание лабораторной работы~2.}
\begin{enumerate}
    \item Для заданных графов (случайно сгенерированных в предыдущей работе)
    реализовать упорядочивание графа по алгоритму Фалкерсона.

    \item Сгенерировать весовую матрицу (положительную или с отрицательными весами --
    выбирает пользователь) и найти кратчайший путь для двух выбранных пользователем
    вершин. В результате выводится не только расстояние, но и сам путь в виде
    последовательности вершин. Реализовать алгоритм Флойда-Уоршалла с выводом матрицы
    расстояний.

    \item Сравнить скорости работы реализованных алгоритмов (по количеству итераций).
\end{enumerate}

\textbf{Задание лабораторной работы~3.}
\begin{enumerate}
    \item На основе сгенерированного ориентированного графа получить случайным образом
    матрицы пропускных способностей и стоимости.

    \item Для полученного графа найти максимальный поток по алгоритму
    Форда--Фалкерсона.

    \item Вычислить заданный поток минимальной стоимости (в качестве величины потока
    брать значение, равное $\lfloor 2/3 \cdot \mathit{max} \rfloor$, где
    $\mathit{max}$~-- максимальный поток). Для поиска кратчайшего пути по стоимости
    использовать алгоритм Флойда-Уоршалла.
\end{enumerate}

\textbf{Задание лабораторной работы~4.}
\begin{enumerate}
    \item Для заданных неориентированных графов (случайно сгенерированных в первой
    работе) найти число остовных деревьев, используя матричную теорему Кирхгофа.

    \item Построить минимальный по весу остов для сгенерированного (неориентированного)
    взвешенного графа, используя алгоритм Краскала. Полученный остов закодировать
    с помощью кода Прюфера и декодировать его; веса сохраняются при кодировании.

    \item Реализовать алгоритм нахождения максимального независимого множества вершин
    на исходном графе и полученном остове (по выбору пользователя).
\end{enumerate}

\textbf{Задание лабораторной работы~5.}
\begin{enumerate}
	\item Для заданных неориентированных графов (случайно сгенерированных в первой работе) проверить, является ли граф эйлеровым. Если нет, то модифицировать граф (логировать, что изменено). Построить эйлеров цикл.
	
	\item Для заданных неориентированных графов (случайно сгенерированных в первой работе) построить кратчайший остов (алгоритмом из четвертой работы).  На основе данного остова получить фундаментальную систему циклов (вывести ее на экран), получение всех остальных циклов на основе фундаментальных с помощью операции симметрической разности (выбранные циклы вводит пользователь).
\end{enumerate}